\documentclass[12pt,a4paper]{article}
\usepackage[utf8]{inputenc}
\usepackage[T1]{fontenc}
\usepackage[italian]{babel}
\usepackage{libertinus}
\usepackage{amssymb}
\usepackage[a4paper,top=25mm,bottom=25mm,left=22mm,right=22mm,headheight=15pt]{geometry}
\usepackage{graphicx}
\usepackage[table]{xcolor}
\usepackage{tabularx,array}
\usepackage{titlesec}
\usepackage{titletoc}
\usepackage{microtype}
\usepackage{fancyhdr}
\usepackage{enumitem}
\usepackage[hidelinks]{hyperref}

% Palette colori del gruppo Bending Forks
\definecolor{bfblue}{HTML}{1B3F9E}
\definecolor{bfrow}{HTML}{F2F5FC}
\definecolor{bfline}{HTML}{C9D1E6}
\definecolor{bfhead}{HTML}{2E6BE6}

\newcommand{\tightlist}{\setlength{\itemsep}{1pt}\setlength{\parskip}{0pt}}
\setlist{topsep=2pt}
\setlist[itemize]{leftmargin=1.3em,label=\textbullet}
\setlist[enumerate]{leftmargin=1.6em}
\setlength{\parindent}{0pt}
\sloppy
\setlength{\parskip}{5pt}
\renewcommand{\arraystretch}{1.45}
\setcounter{secnumdepth}{3}
\setcounter{tocdepth}{3}

% Formattazione titoli sezioni
\titleformat{\section}{\LARGE\bfseries\color{bfblue}}{\thesection.}{0.3em}{}[{\color{bfblue}\vspace{-2pt}\titlerule[0.9pt]}]
\titleformat{\subsection}{\Large\bfseries\color{bfblue}}{\thesubsection.}{0.3em}{}
\titleformat{\subsubsection}{\normalsize\bfseries}{\thesubsubsection.}{0.3em}{}
\titlespacing*{\section}{0pt}{14pt}{6pt}
\titlespacing*{\subsection}{0pt}{10pt}{4pt}
\titlespacing*{\subsubsection}{0pt}{10pt}{3pt}

% Indice cliccabile con puntini/trattini
\contentsmargin{2.2em}
\newcommand{\bfleader}{\titlerule*[0.32em]{\textendash}}
\titlecontents{section}[0pt]{\addvspace{5pt}\large\bfseries}{\thecontentslabel.\ }{}{\hspace{0.4em}\bfleader\contentspage}
\titlecontents{subsection}[1.2em]{\addvspace{0pt}}{\thecontentslabel.\ }{}{\hspace{0.4em}\bfleader\contentspage}
\titlecontents{subsubsection}[2.4em]{\addvspace{0pt}}{\thecontentslabel.\ }{}{\hspace{0.4em}\bfleader\contentspage}

% Intestazioni e piè di pagina
\pagestyle{fancy}
\fancyhf{}
\fancyhead[L]{\footnotesize\sffamily\bfseries\color{bfhead}\textls[120]{BENDING FORKS}}
\fancyhead[R]{\footnotesize\sffamily Verbale Interno N. 2}
\fancyfoot[C]{\small Pagina \thepage}
\renewcommand{\headrulewidth}{0.9pt}

% Colorazione alternata per le tabelle
\newcommand{\bftable}{\rowcolors{2}{bfrow}{white}}
\usepackage{etoolbox}
\AtBeginEnvironment{tabularx}{\bftable}

\begin{document}

% ---------- Copertina ----------
\begin{titlepage}
\centering
\vspace*{4mm}
\includegraphics[width=105mm]{Assets/logo-bending-forks-trasparente.png}\\[4mm]
\underline{bending.forks.swe@gmail.com}

\vspace{42mm}
{\fontsize{24}{28}\selectfont\bfseries Verbale Interno N. 2\par}
\vspace{4mm}
{\Large\color{bfblue} Riunione del 2026-10-19\par}

\vfill
\begin{tabular}[t]{r|l}
\textbf{Responsabile} & Aldo Toffoli \\
\textbf{Scriba} & Luca Stangherlin \\
\textbf{Moderatore} & Filippo Rampazzo \\
\textbf{Verificatore} & Filippo Rampazzo \\
\end{tabular}
\hspace{14mm}
\begin{tabular}[t]{r|l}
\textbf{Versione} & 1.0 \\
\textbf{Tipo} & Verbale Interno \\
\textbf{Destinatari} & \emph{Gruppo Bending Forks} \\
 & \emph{Prof. Tullio Vardanega} \\
 & \emph{Prof. Riccardo Cardin} \\
\end{tabular}

\vspace{12mm}
{\large\emph{\textbf{Abstract dei contenuti:}}}\\[1mm]
\begin{minipage}{150mm}\centering A seguito del primo incontro dove sono stati decisi tre capitolati da approfondire, in vista della candidatura, in questa riunione del 2026-10-19 si è discusso approfonditamente sulla fattibilità e sui dubbi legati a ognuno e in particolare sui quesiti da rivolgere alle aziende. \end{minipage}
\vspace*{6mm}
\end{titlepage}

% ---------- Indice ----------
\vspace*{2mm}
{\fontsize{22}{26}\selectfont\bfseries Indice\par}
\vspace{6mm}
\makeatletter\@starttoc{toc}\makeatother
\newpage

% ---------- Registro delle versioni ----------
\hypertarget{registro-delle-versioni}{%
\section{Registro delle versioni}}

\begin{center}
\begin{tabularx}{\linewidth}{>{\hsize=0.60\hsize\raggedright\arraybackslash}X>{\hsize=0.75\hsize\raggedright\arraybackslash}X>{\hsize=0.95\hsize\raggedright\arraybackslash}X>{\hsize=1.70\hsize\raggedright\arraybackslash}X}
\rowcolor{bfblue}\textcolor{white}{\textbf{Versione}} & \textcolor{white}{\textbf{Data}} & \textcolor{white}{\textbf{Autore}} & \textcolor{white}{\textbf{Descrizione}} \\
1.0 & 2026/10/23 & Luca Stangherlin & Creazione e redazione iniziale del documento \\
1.0 & 2026/10/30 & Filippo Rampazzo & Verifica e approvazione del documento \\
\arrayrulecolor{bfline}\hline
\end{tabularx}
\end{center}

% ---------- Informazioni generali ----------
\hypertarget{informazioni-generali}{%
\section{Informazioni generali}}

\hypertarget{durata-e-luogo}{%
\subsection{Durata e luogo}}

\begin{itemize}
\tightlist
\item \textbf{Data:} 2026-10-19
\item \textbf{Inizio:} 14:30
\item \textbf{Fine:} 16:30
\item \textbf{Luogo / Modalità:} in remoto, tramite Discord
\end{itemize}

\hypertarget{partecipanti}{%
\subsection{Partecipanti}}

\begin{center}
\begin{tabularx}{\linewidth}{>{\raggedright\arraybackslash}X>{\centering\arraybackslash}X>{\centering\arraybackslash}p{2.2cm}>{\centering\arraybackslash}p{2.2cm}}
\rowcolor{bfblue}\textcolor{white}{\textbf{Nome e Cognome}} & \textcolor{white}{\textbf{Ruolo Progetto}} & \textcolor{white}{\textbf{Presente}} & \textcolor{white}{\textbf{Assente}} \\
Marco Beltrame & Membro & $\blacksquare$ & $\square$ \\ \arrayrulecolor{bfline}\hline
Lorenzo Milanese & Membro & $\blacksquare$ & $\square$ \\ \arrayrulecolor{bfline}\hline
Filippo Rampazzo & Verificatore & $\blacksquare$ & $\square$ \\ \arrayrulecolor{bfline}\hline
Nicholas Salmaso & Membro & $\blacksquare$ & $\square$ \\ \arrayrulecolor{bfline}\hline
Luca Stangherlin & Scriba & $\blacksquare$ & $\square$ \\ \arrayrulecolor{bfline}\hline
Aldo Toffoli & Responsabile & $\blacksquare$ & $\square$ \\ \arrayrulecolor{bfline}\hline
Matteo Wu & Membro & $\blacksquare$ & $\square$ \\ \arrayrulecolor{bfline}\hline
\end{tabularx}
\end{center}



% ---------- Ordine del giorno ----------
\hypertarget{ordine-del-giorno}{%
\section{Ordine del giorno}}

\begin{enumerate}
\tightlist
\item {Discussione quesiti su capitolato 1}
\item {Discussione quesiti su capitolato 3}
\item {Discussione quesiti su capitolato 7}
\item {Discussione sulla scelta del capitolato per cui candidarsi}
\end{enumerate}

% ---------- Discussione e decisioni ----------
\newpage
\hypertarget{discussione-e-decisioni}{%
\section{Discussione e decisioni}}

\hypertarget{argomento-1}{%
\subsection{Discussione quesiti su capitolato 1}}

\textbf{Discussione:} \\
{
La discussione su VIMAR si è concentrata sull'integrazione dell'IA e sulla sicurezza domotica. Si è analizzato come l'utente possa collegare l'LLM al server MCP su Web, desktop e mobile, e quali componenti fornire per abilitare la comunicazione. L'utilizzo dell'autenticazione VIMAR, necessaria per verificare l'identità dell'utente, è stato un punto di discussione e il centro delle domande.
 }

\textbf{Decisione presa:}
\begin{itemize}
\tightlist
\item[$\checkmark$] \textbf {Viene presa la decisione di inoltrare a Vimar il form per un colloquio da remoto.}
\end{itemize}

\textbf{Motivazione:} \\
{Sono state scelte le domande che lasciavano reali lacune per una profonda analisi del progetto.}

\textbf{Responsabile attuazione:} Filippo Rampazzo \\
\textbf{Scadenza:} 2026-10-19

\hypertarget{argomento-2}{%
\subsection{Discussione quesiti su capitolato 3}}

\textbf{Discussione:} \\
{
Ci siamo confrontati sulla modalità delle sessioni di formazione e sugli aspetti hardware e audio, in particolare per chiarire la gestione di eco, beamforming e dispositivi multipli. Abbiamo inoltre trattato la scelta dell'infrastruttura AI, confrontando  crediti Cloud o VPS locale e definito come gestire l'estensione del sistema tramite skill. Si è discusso di una dashboard centralizzata per utenti e abitazioni.
}

\textbf{Decisione presa:}
\begin{itemize}
\tightlist
\item[$\checkmark$] \textbf {Viene presa la decisione di inoltrare a Uranion una mail per un colloquio da remoto.}
\end{itemize}

\textbf{Motivazione:} \\
{Sono state scelte le domande che lasciavano reali lacune per una profonda analisi del progetto.}

\textbf{Responsabile attuazione:} Aldo Toffoli\\
\textbf{Scadenza:} 2026-10-19

\hypertarget{argomento-3}{%
\subsection{Discussione quesiti su capitolato 7}}

\textbf{Discussione:} \\
{
Le domande per l'incontro con Ergon hanno riguardato i temi sulla personalizzazione della Smart Dashboard, sul controllo di eventuali allucinazioni degli LLM, sull'uso di quest'ultimi in modalità gratuita piuttosto che a pagamento e infine sulle modalità e le periodicità dei confronti con l'azienda.
}

\textbf{Decisione presa:}
\begin{itemize}
\tightlist
\item[$\checkmark$] \textbf {Viene presa la decisione di inoltrare a  una mail a Ergon per un colloquio da remoto.}
\end{itemize}

\textbf{Motivazione:} \\
{Sono state scelte le domande che lasciavano reali lacune per una profonda analisi del progetto.}

\textbf{Responsabile attuazione:} Filippo Rampazzo\\
\textbf{Scadenza:} 2026-10-19

\hypertarget{argomento-4}{%
\subsection{Discussione sulla scelta del capitolato per cui candidarsi}}

\textbf{Discussione:} {
Si è svolto un dibattito su quale capitolato dei tre selezionati potesse riflettere la scelta migliore per il team. Vi è stata una divisione di quest'ultimo tra i sostenitori del capitolato 1 e i sostenitori del capitolato 3. Riconoscendo  i pro e i contro di entrambi si è arrivati ad una decisione. 
}

\textbf{Decisione presa:} {Il team sceglie di portare avanti la canditatura per il capitolato 3}


\textbf{Motivazione:} Le discussioni interne seguite da relative votazioni dei capitolati hanno portato alla scelta del capitolato 3.

\textbf{Responsabile attuazione:} Team Bending Forks\\
\textbf{Scadenza:} /



% ---------- Azioni da intraprendere (TODO) ----------
\hypertarget{azioni-da-intraprendere}{%
\section{Azioni da intraprendere (TODO)}}

\begin{center}
\begin{tabularx}{\linewidth}{>{\hsize=1.5\hsize\raggedright\arraybackslash}X>{\hsize=1.0\hsize\raggedright\arraybackslash}X>{\hsize=0.8\hsize\centering\arraybackslash}X>{\hsize=0.7\hsize\centering\arraybackslash}X}
\rowcolor{bfblue}\textcolor{white}{\textbf{Descrizione}} & \textcolor{white}{\textbf{Responsabile}} & \textcolor{white}{\textbf{Scadenza}} & \textcolor{white}{\textbf{Priorità}} \\
Invio form a Vimar & Filippo Rampazzo & 2026-10-19 & Alta \\ \arrayrulecolor{bfline}\hline
Invio mail a Uranion & Aldo Toffoli & 2026-10-19 & Alta \\ \arrayrulecolor{bfline}\hline
Invio mail a Ergon & Filippo Rampazzo & 2026-10-19 & Alta \\ \arrayrulecolor{bfline}\hline
\end{tabularx}
\end{center}

% ---------- Prossima riunione ----------
\hypertarget{prossima-riunione}{%
\section{Pianificazione prossima riunione}}

\begin{itemize}
\tightlist
\item \textbf{Data proposta:} 2026-10-23
\item \textbf{Ora prevista:} 16:30
\item \textbf{Luogo / Canale:} Canale Discord
\item \textbf{OdG preliminare:}
    \begin{itemize}
    \item[-] Comprensione documenti 30 ottobre
    \item[-] Repository GitHub
    \item[-] Definizione del way of working
    \item[-] Organizzazione struttura diario di bordo
    \end{itemize}
\end{itemize}


\vspace{35pt}
\noindent
\begin{minipage}[t]{0.45\textwidth}
\textbf{Firma Scriba:}\\[18pt]
\makebox[0pt][l]{\raisebox{0mm}{\includegraphics[width=40mm]{Assets/firma_luca.png}}}\rule{55mm}{0.5pt}\\[2pt]
\small Luca Stangherlin
\end{minipage}
\hfill
\begin{minipage}[t]{0.45\textwidth}
\begin{flushright}
\textbf{Firma Verificatore:}\\[18pt]
\makebox[0pt][l]{\raisebox{0mm}{\includegraphics[width=40mm]{Assets/firma_filippo.jpeg}}}\rule{55mm}{0.5pt}\\[2pt]
\small Filippo Rampazzo
\end{flushright}
\end{minipage}

\end{document}